\section{Further attachment and subdivision consequences}\label{eight:extensions}

\begin{corollary}\label{eight:common-short}
Let $q\ge1$, $u\ge0$, and $0\le r\le8$.  If $b_1,\ldots,b_r\ge3$
are odd, then
\[
 \Sp{qb_1,\ldots,qb_r,q^u}
\]
is graceful.  Likewise, under the hypotheses of
Article Theorem~\ref*{main-art:even-eight},
$\Sp{qa,qb_1,\ldots,qb_7,q^u}$ is graceful.
In particular every eight-arm spider whose lengths have the same
$2$-adic valuation is graceful.
\end{corollary}
\begin{proof}
Apply the respective theorem before subdivision, with $u$ unit arms,
and then use ordinary full $q$-subdivision
(Theorem~\ref{common:offset-subdivision}, at source offset zero).
Root the ordinary source labeling at whichever vertex has label zero;
that vertex need not be the hub.  The underlying subdivided spider has
the displayed lengths.  For the final assertion divide all lengths by
their common power of two to obtain eight positive odd quotients.
\end{proof}

\begin{proposition}[Four-path attachment]\label{eight:four-paths}
Let $G$ be any root-zero alpha-labeled tree of index $K$.  Adjoin four
positive-length paths $a,b,x,y$ at its root.  If $a+b=K+2$, the
resulting tree is graceful, also with any number of additional unit
paths at that root.  When $(a,b)\notin\{(2,2),(3,3)\}$, the two
paths $x,y$ may instead be replaced by any root-zero alpha-labeled
tree $T$.
\end{proposition}
\begin{proof}
Suppose first that $(a,b)\notin\{(2,2),(3,3)\}$.  The maximal-pair
theorem in Section~\ref{app:pair-offsets} gives offset $a+b-2=K$ on
$\Sp{a,b}$.  Alpha-offset amalgamation with $G$ gives a hub-zero
component.  If $(x,y)\ne(2,2)$, the two further arms form a root-zero
alpha pair, which can be ordinarily amalgamated with that component.
Append the units to the hub-zero component first.

If $(x,y)=(2,2)$, the four attached paths have a short arm, so their
whole four-arm spider is hub-zero.  Ordinarily amalgamate it with $G$,
after adjoining the units to the hub-zero piece.  The same direct
short-arm argument applies if $(a,b)$ is either maximal-offset
exception: in those cases the four-arm piece has an arm at most three.
Thus neither a nonexistent maximal offset nor the alpha-pair exception
is used.  The assertion with a general $T$ follows from the first
hub-zero component and ordinary alpha amalgamation; its stated
restriction on $(a,b)$ is retained.
\end{proof}

\begin{corollary}\label{eight:equal-four-six-core}
For every odd $L\ge7$ and all positive $a,b$ with $a+b=2L$,
$\Sp{L,L,L,L,a,b}$ is hub-zero graceful.  Adjoining arbitrary positive
arms $x,y$ and any number of unit arms gives a graceful spider.
\end{corollary}
\begin{proof}
The equal-odd-four alpha component has index $2L-2$, and the pair has
maximal offset $a+b-2=2L-2$.  The two forbidden maximal pairs have
sums four and six, below $2L$, so neither occurs.  This proves the
hub-zero claim.  The ordinary extension also follows directly from
Proposition~\ref{eight:four-paths}.
\end{proof}

\begin{proposition}[Five long paths]\label{eight:five-long-paths}
Let $G$ be any root-zero alpha-labeled tree of index $K$.  Attach five
paths with lengths $L_1,\ldots,L_5$ satisfying
\[
 L_i>2K,\qquad L_i\ne4K+1\quad(1\le i\le5).
\]
The resulting tree is graceful, also after adjoining arbitrary unit
paths at the same root.
\end{proposition}
\begin{proof}
Select a hub-zero four-arm piece among these five paths by
Section~\ref{app:five-to-four}.  The unused path attaches to $G$ as a
root-zero alpha extension by Lemma~\ref{common:alpha-extension}.
Ordinarily amalgamate it with the selected hub-zero piece, after
appending the units to that piece.  The hypotheses refer only to the
original index $K$; no successive growth or hierarchy among the five
lengths is needed.
\end{proof}

\begin{proposition}[A mixed five-arm core]\label{eight:double-five-core}
For every $r\ge2$, the spider
\[
 F_r=\Sp{2r,4r,2r+1,2r+1,2r+1}
\]
is hub-zero graceful.  Ordinarily amalgamating $F_r$ with any
root-zero alpha tree $T$, and adjoining any number of unit arms, gives
a graceful tree.  In particular, for arbitrary odd $a,b,c\ge5$,
\[
 \Sp{2r,4r,(2r+1)^3,a,b,c,1^u}
\]
is graceful for every $u\ge0$.
\end{proposition}
\begin{proof}
The literal pair $(2)\ ;\ (3,4)$ has offset one on $\Sp{1,2}$.
Subdivide by $2r$ and choose residue $r$.  The new pair has lengths
$2r,4r$ and offset $2r+r=3r$.  The equal odd alpha triple has common
length $2r+1\ge5$ and index $3r$.  Alpha-offset amalgamation yields
$F_r$.  Append the units there before ordinary amalgamation with $T$.
For the last statement use the universal odd alpha triple $\Sp{a,b,c}$.
The restriction $r\ge2$ avoids the unavailable source alpha triple
$\Sp{3,3,3}$.

For completeness the offset pair needs no path search.  Form $A_r$
by concatenating the pairs $(4r-j,8r+j+1)$ in the order
$j=0,\ldots,r-1$.  Form $B_r$ by first concatenating all pairs
$(6r-j,6r+j+1)$ in that order, and then concatenating all pairs
$(5r-j,7r+j+1)$ in that order.
These use $3r+1,\ldots,9r$.  The internal colors of $A_r$ are
$4r+1,\ldots,6r-1$ and its root color is $4r$.  The internal colors
of $B_r$ are $1,\ldots,4r-1$ and its root color is $6r$.
\end{proof}

\begin{proposition}[Three equal paths]\label{eight:three-equal-offsets}
For every $q\ge1$,
\[
 [0,7q-1]\subseteq\Om(\Sp{3q,3q,3q}).
\]
Consequently, if a root-zero alpha tree $G$ has index $K\le7q-1$,
attaching three paths of length $3q$ gives a hub-zero graceful tree.
Ordinarily amalgamating any second root-zero alpha tree $T$ gives a
graceful tree, also with arbitrary unit arms.
\end{proposition}
\begin{proof}
Here are offset-$s$ certificates on $\Sp{3,3,3}$ for each $s=0,\ldots,6$:
\[
\begin{array}{c|ccc}
s&\multicolumn{3}{c}{\text{words outward from root zero}}\\ \hline
0&(4,7,5)&(8,1,6)&(9,3,2)\\
1&(7,3,6)&(8,2,4)&(9,10,5)\\
2&(6,5,10)&(7,4,8)&(9,11,3)\\
3&(6,4,12)&(7,11,8)&(9,10,5)\\
4&(5,12,6)&(8,7,10)&(9,13,11)\\
5&(7,13,14)&(8,12,10)&(9,6,11)\\
6&(7,13,15)&(8,11,12)&(9,14,10)
\end{array}
\]
Each row has labels $s+1,\ldots,s+9$ and colors $1,\ldots,9$.
Factor-$q$ offset subdivision supplies every offset in
$[qs,q(s+1)-1]$.  The seven intervals partition $[0,7q-1]$.
The attachment and ordinary-gluing conclusions follow from the common
amalgamation lemmas.  Append units to a hub-zero component first.
\end{proof}

\begin{corollary}\label{eight:nine-family}
Suppose a nine-arm spider has three equal arms of length $3q$, with
$q\ge1$, and its other six lengths are odd integers at least five.
If three of those six have half-length sum at most $7q-1$, the spider
is graceful, also with arbitrary unit arms.  The remaining three odd
lengths have no upper bound.  For odd $q$ all nine arms are odd; for
even $q$ the three designated arms are even.
\end{corollary}
\begin{proof}
Use the bounded-index odd triple as $G$ in
Proposition~\ref{eight:three-equal-offsets} and the other odd triple as
$T$.  Both are root-zero alpha since their minima are at least five.
\end{proof}

\begin{proposition}[Upper pair bands]\label{eight:dyadic-pair-band}
Let $a,b$ be positive and let $d=2^h$ divide both, where $h\ge0$.
If $a\ne b$, then
\[
 [a+b-2d,a+b-2]\subseteq\Om(\Sp{a,b}).
\]
The same conclusion holds for $a=b$ if $a\ge4d$.
\end{proposition}
\begin{proof}
Write $m=a+b$.  For $q=1,2,4,\ldots,d$, the quotient pair is unequal
and hence avoids $(2,2),(3,3)$.  Its maximal offset is $m/q-2$.
Factor-$q$ subdivision gives
\[
 [m-2q,m-q-1].
\]
These intervals abut and their union is $[m-2d,m-2]$.
For equal arms the smallest quotient is $a/d\ge4$, so every quotient
still avoids both forbidden maximal pairs and the same proof applies.
\end{proof}

Uniform full subdivision may be applied to each ordinary graceful
family proved in this section.  This multiplies every arm length by
the subdivision factor and turns each unit arm into an arm of that
common length.  It uses ordinary subdivision, with no assertion that
alpha separation is preserved.
